\cotefile{132}
\scanpage{2}
\mspage{1}
\diagram{\begin{tikzcd}[ampersand replacement=\&,column sep=small,row sep=small]
 X \arrow[dr,"f"'] \& \& Y \arrow[dl,"g"]\\
 \& S \&
\end{tikzcd}}
espaces topologiques, d'où $X\times_S Y$

et
\[
\Pi_1(X\times_S Y)\xrightarrow{\varphi}
\Pi_1(X)\mathbin{\overset{2}{\times}}_{\Pi_1(S)}\Pi_1(Y)
\]

a) Supposons $f$ une fibration de Serre, alors $\varphi$ induit une bijection sur les $\pi_0$, un épimorphisme sur les $\pi_1$ associés aux composantes connexes (i.e. $\varphi$ \deleted{est} ess. surjectif et surjectif sur les $\mathrm{Hom}(x,y)\to\mathrm{Hom}\,\varphi x,\varphi y$).

\uline{Démonstration} Découpant la situation \deleted{\unsure{suivant}} les composantes connexes de $X,S,Y$, on est ramené au cas où $S,X,Y$ sont connexes\inserted{ non vides}, \deleted{\unsure{donc|alors}} donc $f$ surjectif. Comme $Y\ne\varnothing$, prenant $y_0\in Y$, $x_0\in X_{s_0}$ où $s_0=g(y_0)$. Posons
\[
G=\pi_1(X,x_0),\quad H=\pi_1(Y,y_0),\quad \Sigma=\pi_1(S,s_0)
\]
on a alors
\diagram{\begin{tikzcd}[ampersand replacement=\&,column sep=small,row sep=small]
 G \arrow[dr,"f_1"'] \& \& H \arrow[dl,"g_1"]\\
 \& \Sigma \&
\end{tikzcd}}
permettant de reconstruire mod équivalence
\[
\Pi_1(X)\mathbin{\overset{\unsure{2|?}}{\times}}_{\Pi_1(S)}\Pi_1(Y)
\quad\text{comme}\quad
\mathrm{Cat}\,G\mathbin{\overset{2}{\times}}_{\mathrm{Cat}\,\Sigma}\mathrm{Cat}\,H=:P
\]
Les objets de cette dernière sont les objets $\sigma$ de $\Sigma$, et les homs $\sigma\to\sigma'$ sont les couples $(g,h)\in G\times H$, avec $g_1(h)\sigma=\sigma'f_1(g)$, i.e. $\sigma'=g_1(h)\sigma f_1(g)^{-1}$. Donc les classes d'isomorphie de $P$ correspondent à l'ensemble des doubles classes $H\backslash\Sigma/G$, et \deleted{les $\pi_1$ correspondants} \inserted{$\pi_1(P,\sigma)$ n'est autre} \deleted{aux \unsure{groupes|composantes} des \unsure{auto}} \inserted{que le stabilisateur} de $\sigma\in\Sigma$ dans $G\times H$, opérant sur $\Sigma$ par $(h,g)\cdot\sigma=g_1(h)\sigma f_1(g)^{-1}$
\margin{\diagram{\begin{tikzcd}[ampersand replacement=\&,column sep=small,row sep=tiny]
e_G \& e_H\\
e_\Sigma \arrow[r,"\sigma"] \arrow[d,"f_1(g)"'] \& e_\Sigma \arrow[d,"g_1(h)"]\\
e_\Sigma \arrow[r,"\sigma'"'] \& e_\Sigma
\end{tikzcd}}}
% archive sheet number: 1
\scanpage{3}
\mspage{2}
Calculons $\pi_0(Z)$ $(Z=X\times_S Y)$, en utilisant la fibration de Serre
\[
 F\hookrightarrow Z\longrightarrow Y
\]
où $F$ est la fibre de $f:X\to S$, on trouve \inserted{(comme $\pi_0(Y)=1$)} (avec les points base déjà choisis, \uline{et} $z_0=(x_0,y_0)$)
\[
 \pi_0(Z)=\pi_0(F)/\pi_1(Y)\ \text{(opérant à g.)}
\]
d'ailleurs on a, par la fibration de Serre (et $\pi_0(X)=1$)
\[
 F\hookrightarrow X\longrightarrow S
\]
\[
 \deleted{\pi_1(X)\to\pi_1(S)\to\pi_0(F)\to\pi_0\unsure{(X)}\to1}
\]
\deleted{i.e.}
\[
 \pi_0(F)\simeq\pi_1(S)/\underbrace{\pi_1(X)}_{\text{op. à droite}}
\]
d'où, en composant
\[
 \pi_0(Z)=\pi_1(Y)\backslash\pi_1(S)/\pi_1(X)
\]
ce qui montre la bijectivité des $\pi_0$.

\unsure{Récipt.|Reste.} Soit $Z'$ une composante connexe de $Z$. Elle se projette \uline{sur} $Y$, donc $\exists\,z'_0\in Z'_{y_0}$. Son image dans $X$ est \deleted{bien} dans $X_{s_0}=F$, donc
% slanted
\margin{donner une description symétrique en $X,Y$ \ldots}
en la reliant par un arc à $x_0$, on trouve en appliquant $f$ un lacet en $s_0$, soit $\sigma$ [déterminé mod un élément de \inserted{$\mathrm{Im}$} $\pi_1(X,x_0)=G$ par chgt. de l'arc dans $X\ldots$, et modulo un élément de $\mathrm{Im}\,\pi_1(Y,y_0)$ en changeant le point $z'_0$ \ldots.]
% archive sheet number: 2
\scanpage{4}
\mspage{3}
\diagram{\scriptsize\begin{tikzcd}[ampersand replacement=\&,column sep=.55em,row sep=1.3em]
{[\pi_2(Z)} \arrow[r] \& {\pi_2(Y)]} \arrow[r] \arrow[d] \& \pi_1(F) \arrow[r] \arrow[d,"\text{épi}"] \& \underline{\pi_1(Z)} \arrow[r] \arrow[d] \& \pi_1(Y) \arrow[r] \arrow[d,"\unsure{\sim|2|\wr}"] \& \pi_0(F) \arrow[r] \arrow[d,"\sim"] \& \underline{\pi_0(Z)} \arrow[r] \arrow[d] \& 1\\
\& \deleted{\unsure{*}}\ 1 \arrow[r] \& \pi_1(\Phi) \arrow[r] \& \pi_1(\widetilde Z) \arrow[r] \& \pi_1(\widetilde Y) \arrow[r] \arrow[d,phantom,"\wr"] \& \pi_0(\Phi) \arrow[r] \arrow[d,phantom,"\wr"] \& \pi_0(\widetilde Z) \arrow[r] \& 1\\
\& \& \& \& H \& \Sigma/G
\end{tikzcd}}
\diagram{\begin{tikzcd}[ampersand replacement=\&,row sep=tiny,column sep=1.3em]
F \arrow[d,hook] \& \deleted{\unsure{\Phi}} \& \Phi \arrow[d] \&\\
X \arrow[d] \& \deleted{\pi_1(X)=G} \arrow[d] \& \mathrm{Cat}\,G \arrow[d] \arrow[r,phantom,"\simeq"] \& \widetilde X \arrow[d]\\
S \& \deleted{\pi_1(S)} \& \mathrm{Cat}\,\Sigma \arrow[r,phantom,"="] \& \widetilde S
\end{tikzcd}}
$\Phi={}$ Catégorie des \deleted{éléments} de $\Sigma$, avec $\mathrm{Hom}_{\Phi}(\sigma,\sigma')={}$ ens des $g\in G$ tels que $\sigma=\sigma'f_1(g)$ i.e. $g\in f_1^{-1}(\sigma'^{-1}\sigma)$

Son $\pi_0$ est $\Sigma/G$, ses $\pi_1$ tous isom. à $\mathrm{Ker}\,f_1$.

On a $\Pi_1(F)\longrightarrow\Phi$, induisant une bijection sur les $\pi_0$, et sur les $\pi_1$ un épim., dont le noyau est $\mathrm{Coker}\,(\pi_2(X)\to\pi_2(S))$
\diagram{\scriptsize\begin{tikzcd}[ampersand replacement=\&,row sep=1.1em,column sep=.7em]
\pi_2(X) \arrow[r] \& \pi_2(S) \arrow[r] \& \pi_1(F) \arrow[r] \& \pi_1(X) \arrow[r] \arrow[d,phantom,"\wr"] \& \pi_1(S) \arrow[r] \arrow[d,phantom,"\wr"] \& \pi_0(F) \arrow[r] \arrow[d,phantom,"\wr"] \& 1\,\deleted{\unsure{\longrightarrow\pi_0}}\\
\& 1 \arrow[r] \& \pi_1(\widetilde F) \arrow[r] \& \pi_1(\widetilde X) \arrow[r] \& \pi_1(\widetilde S) \arrow[r] \& \pi_0(\Phi) \arrow[r] \& 1
\end{tikzcd}}

\[
\sigma=1\qquad \pi_1(\widetilde Z)=\{(g,h)\mid f_1(g)=g_1(h)
\]
$\pi_1(\widetilde Z)$ est ext. de $g_1^{-1}(\mathrm{Im}\,G)$ par $\pi_1(\Phi)$

$\pi_1(Z)$ est ext. de idem par $\mathrm{Coker}\,(\pi_2(Y)\to\pi_1(F))$

donc est ext. de $\pi_1(\widetilde Z)$ par
\[
\mathrm{Coker}\,(\pi_2(X)\times\pi_2(Y)\to\pi_2(S))
\]

\uline{Donc on trouve}

$Z\longrightarrow\widetilde Z$ induit un iso. pour les $\pi_1$ sss $\pi_2(X)\times\pi_2(Y)\longrightarrow\pi_2(S)$ est épi.

De plus, $\pi_i(Z)=0$ pour $i\geq 2$ \deleted{\unsure{si|\Leftrightarrow}} dès que $\pi_i(Y)=\pi_i(F)=0$ pour $i\unsure{=|\geq}2$, or $\pi_{i+1}(S)\to\pi_i(F)\to\pi_i(X)$ donc il suffit $\pi_i(X)=\pi_i(Y)=\pi_i(S)=0$ pour $i\geq 2$, i.e. $X,Y,S$ sont des $1$-tronqués\,\ldots
% archive sheet number: 3
\scanpage{6}
\diagram{\begin{tikzcd}[ampersand replacement=\&,column sep=large]
\underset{\substack{\text{groupoïde}\\\text{1-spécial}\\\text{orienté}}}{D}
\arrow[r,"\varphi", "\substack{\text{foncteur}\\\text{spécial}}"'] \&
\underset{\substack{\text{groupoïde}\\\text{2-spécial}}}{C}
\end{tikzcd}}
$0$-objets \qquad $(C,D,\varphi:D\to C)$

\diagram{\begin{tikzcd}[ampersand replacement=\&]
D \arrow[r,"\varphi"] \arrow[d,"u"',"\wr"] \& C \arrow[d,"\deleted{\unsure{u}}\,v"]\\
D' \arrow[r,"\varphi'"'] \& C' \arrow[ul,phantom,very near start,"\alpha\,\unsure{\swarrow}"]
\end{tikzcd}}
$1$ objets \qquad $(u,\underset{\substack{\downarrow\\\text{compatible}\\\text{avec orientations}}}{v},\alpha:\varphi'u\xrightarrow{\sim}v\varphi)$

\diagram{\begin{tikzcd}[ampersand replacement=\&,column sep=large]
D \arrow[r,"\varphi"] \arrow[d,"u_1"',""{name=U,coordinate,near start}] \arrow[d,bend left=48,"u_2",""{name=V,coordinate,near start}] \&
C \arrow[d,"v_1"',""{name=W,coordinate,near start}] \arrow[d,bend left=48,"v_2",""{name=T,coordinate,near start}]\\
D' \arrow[r,"\varphi'"'] \& C' \arrow[ul,phantom,near start,"\alpha_1\Rightarrow\quad\alpha_2\Rightarrow"]
\arrow[from=U,to=V,Rightarrow,"\unsure{\lambda}"] \arrow[from=W,to=T,Rightarrow,"\mu"]
\end{tikzcd}}
$2$-objets \qquad $(\lambda,\mu)$ \qquad
$\begin{aligned}\lambda &:u_1\overset{\sim}{\longrightarrow}u_2\\
\mu &:v_1\overset{\sim}{\longrightarrow}v_2\end{aligned}$

\diagram{\begin{tikzcd}[ampersand replacement=\&]
\varphi'u_1 \arrow[r,"\varphi'*\lambda","\sim"'] \arrow[d,"\alpha_1"',"\wr"] \&
\varphi'u_2 \arrow[d,"\alpha_2","\wr"']\\
\deleted{\unsure{\varphi}}v_1*\varphi \arrow[r,"\mu*\varphi","\sim"'] \& v_2*\varphi
\end{tikzcd}\quad commutatif}

supposons $C$ connexe \inserted{non vide}, \inserted{$D$ non vide}, alors $\varphi:D\to C$ est équivalent à un objet (noté de $\overline{\rm m}$) où
\[
 C=\mathrm{Cat}\,G,\qquad D=\coprod_{i\in I}\mathrm{Cat}\,\mathbb Z
\]
donc $\varphi$ donné par des homs $\varphi_i:\mathbb Z\to G$

donc par des $h_i=\varphi_i(1)\in G$ \qquad $(i\in I$

idem \qquad $C'=\mathrm{Cat}\,G',\quad D'=\coprod_{i'\in I'}\mathrm{Cat}\,\mathbb Z$

$\unsure{\psi|\varphi|\varphi'}$ donné par des $\varphi'_{i'}:\mathbb Z\to G'$ i.e. des $h'_{i'}\in G'$

\[
\underline{\mathrm{Hom}}\big((C,D,\varphi);(C',D',\varphi')\big)\deleted{\unsure{=}}\overset{?}{=}
\]
\[
\text{Objets}\ \left\{\begin{aligned}
&\boxed{\begin{aligned}u &:I\xrightarrow{\sim}I'\\v &:G\xrightarrow{\sim}G'\end{aligned}}\ 
  \begin{aligned}&\text{(définit $u$)}\\ &\text{(définit $v$)}\end{aligned}\\
&\forall i\in I \quad \boxed{\alpha(i)\in G'}\ \text{tel que}\quad
\deleted{\unsure{\varphi'u}\ \unsure{h'_{u(i)}=\alpha_i v(h_i)\alpha_i^{-1}|\ell'_{u(i)}=\alpha_i v(\ell_i)\alpha_i^{-1}}}\\
&\text{\unsure{(de|(idem}}\hspace{3em}\text{(définit $\alpha$)}\quad
\boxed{v(h_i)=\alpha(i)h'_{u(i)}\alpha(i)^{-1}}
\end{aligned}\right.
\]
\[
\begin{aligned}
&\mathrm{Hom}\big((u_1,v_1,\alpha_1),(u_2,v_2,\alpha_2)\big)
 \underset{\substack{\text{si }u_1=u_2=u\\\text{sinon c'est vide}}}{=}
 \big\{(\lambda,\mu)\mid\lambda\in\mathbb Z^I,\ \mu\in G',\ \deleted{\unsure{\in}}\\
&\qquad \left\{\begin{aligned}
&\boxed{v_2=\deleted{\unsure{\mu}}\mathrm{int}(\mu)\circ v_1}\\
&\alpha_2(i){h'_{u(i)}}^{\lambda_i}=\mu\alpha_1(i)\qquad \text{ou aussi}\\
&\boxed{\alpha_2(i)^{-1}\mu\alpha_1(i)={h'_{u(i)}}^{\lambda_i}}
\end{aligned}\right.
\end{aligned}
\]
% archive sheet number: 5
\scanpage{7}
Supposons $(C',D',\varphi')=(C,D,\varphi)$ donc $G=G'$, $I=I'$, $h'_i=h_i$

donc $\underline{\mathrm{Hom}}(,)=\underline{\mathrm{End}}((C,D,\varphi)$ est une Gr-catégorie.

\[
\pi_0\big(\underline{\mathrm{End}}(C,D,\varphi)\big)\overset{?}{=}
\]

\uline{NB} $(u_1,v_1,\alpha_1)\sim(u_2,v_2,\alpha_2)$ \deleted{\unsure{ssi}} implique

a) $u_1=u_2\ (=u)$

b) $v_2$ conjugué de $v_1$ i.e. $\exists\mu\in G'$ tel que $v_2=\mathrm{int}(\mu)v_1$

\[
\begin{aligned}
\text{d'où}\quad v_2(h_i)&=\mathrm{int}(\mu)\,v_1(h_i)=\mu\,v_1(h_i)\mu^{-1}\\[-2pt]
\Vert\quad &\hspace{11em}\Vert\\[-2pt]
\alpha_2(i)h_{u(i)}^{\deleted{\unsure{\prime}}}\alpha_2(i)^{-1}
&=\mu\alpha_1(i)h_{\unsure{u_1(i)|u_2(i)|u(i)}}^{\deleted{\unsure{\prime}}}\alpha_1(i)^{-1}\mu^{-1}
\end{aligned}
\]
d'où \quad $(\mathrm{int}\,\alpha_2(i))h_{u(i)}=(\mathrm{int}(\mu\alpha_1(i)))\cdot h_{u(i)}$

i.e \deleted{int} $\alpha_2(i)^{-1}\mu\alpha_1(i)^{\deleted{\unsure{-1}}}$ centralise $h_{u(i)}$

Or sauf le cas de la ``sphère à 1 seul trou'', chaque $h_j$ est un élément libre, et \inserted{en tous cas,} son centralisateur est le groupe engendré par $h_j$ ; \deleted{donc, dans les cas, \unsure{il}} donc a) et b) impliquent déjà l'existence des $\lambda_i$ tels que $\alpha_2(i)^{-1}\mu\alpha_1(i)=h_{u(i)}^{\lambda_i}$ (d'ailleurs uniques sauf le cas de la sphère \unsure{à} un seul trou i.e. $G=1$, $\mathrm{card}\ I=1$), donc
\[
(u_1,v_1,\alpha_1)=(u_2,v_2,\alpha_2)\quad\text{sss}\quad u_1=u_2,\ \widetilde v_2=\widetilde v_1
\]
($\sim$ dénote passage au groupe extérieur). Donc

$\pi_0(\underline{\mathrm{End}})\simeq$ groupe des couples $(u,\widetilde v)$, où $u\in\unsure{\mathfrak S_I|\mathfrak G_I}$, $\widetilde v\in\mathrm{Aut\,ext}\,G$, tels que $\forall i\in I$, on ait $\widetilde v(h_i)$ conjugué à $h_{u(i)}$ (pour $v$ un \unsure{auto|aut.} de $G$ qui représente $\widetilde v$).

\uline{NB} $\sigma$ est déjà déterminé par $\widetilde v$, ce qui résulte du fait que si $i,j\in I$, $i\ne j$ alors $h_i$ et $h_j$ sont non conjugués (en fait leurs images dans $G_{\rm ab}$ sont distinctes !) Donc $\pi_0(\underline{\mathrm{End}})\subset\mathrm{Aut\,ext}(G)$, est formé des $\widetilde u\in\mathrm{Aut\,ext}$ tels que $\forall i\in I$, $\widetilde v(h_i)$ est conjugué à un des $h_j$.
% archive sheet number: 6
\scanpage{8}
\[
\pi_1(\underline{\mathrm{End}})\overset{?}{=}
\mathrm{Aut}\,(u=\mathrm{id},v=\mathrm{id},\alpha\deleted{\unsure{\in}}=1\in G^I)=\ ?
\]
$\lambda\in\mathbb Z^I$, $\mu\in G$ tels que

a) $\mathrm{id}_G=\mathrm{int}(\mu)\circ\mathrm{id}_G$ i.e. $\mu\in Z(G)$ (centre)

b) $\forall i,\quad\mu=h_i^{\lambda_i}$

Or, sauf le cas de la sphère à $1$ ou $2$ trous, le groupe $G$ est libre à $\unsure{N|n}\geq2$ générateurs, donc son centre est nul, de plus les $h_i$ sont libres. Donc dans ce cas on a $\mu=1$, $\lambda_i=0\ \forall i\in I$, donc $\boxed{\pi_1(\underline{\mathrm{End}})=0}$

Dans le cas de la sphère à $2$ trous, \inserted{i.e. le cylindre,} on a $G\simeq\mathbb Z$, $h_i=\pm1$, $\deleted{\unsure{h_j=\pm1}}$ ; les $h_i$ sont encore libres (donc les $\lambda_i$ déterminés par $\mu$) et $G=Z(G)$. Donc
\[
\pi_1(\underline{\mathrm{End}})=G\quad(\simeq\mathbb Z)
\]
\margin{(pas canoniquement)}
Dans le cas de la sphère à $1$ trou, $G=0$, et on trouve $\pi_1(\underline{\mathrm{End}})\simeq\mathbb Z$ (isom. canonique)

\begin{center}\rule{1.5cm}{.4pt}\end{center}

Regardons maintenant la $2$-catégorie des $(C,D,\varphi)$ \uline{avec $D$ fixé}, donc $u=\mathrm{id}$, $\lambda=\mathrm{id}$. On trouve pour $\underline{\mathrm{End}}_{D\text{ fixé}}(C,D,\varphi)$

objets : couples $(v,\alpha)$ avec
\[
\left\{\begin{aligned}
&v\in\mathrm{Aut}\,G,\\
&\alpha\in G^I\\
&v(h_i)=\mathrm{int}(\alpha(i))\cdot h_i
\end{aligned}\right.
\]
$\mathrm{Hom}\big((v_1,\alpha_1),(v_2,\alpha_2)\big)={}$ ens. des $\mu\in G$ tels que
\[
\left\{\begin{aligned}
&v_2=\mathrm{int}(\mu)\circ v_1\\
&\alpha_2(i)=\mu\alpha_1(i)
\end{aligned}\right.
\]
On trouve que $(v_1,\alpha_1)\simeq(v_2,\alpha_2)$ $\deleted{\unsure{\sim}}\Rightarrow\deleted{\unsure{\exists}}$ \deleted{\unsure{$\mu$ tel q}} \inserted{$\widetilde v_1=\widetilde v_2$} i.e $\exists\mu_0$ tel que $v_2=\mathrm{int}(\mu_0)v_1$
\[
\begin{aligned}
\text{d'où}\qquad v_2(\unsure{h^i|h_i|h'_i|\ell_i})&=\underline{\mathrm{int}(\mu_0)}\,\underline{v_1}(h_i)\\
\Vert\qquad &\hspace{5em}\Vert\\
\mathrm{int}(\alpha_2(i))h_i &\qquad\mathrm{int}(\mu_0)\,\mathrm{int}\,\alpha_1(i)\,v_1(h_i)
\end{aligned}
\]
donc $\alpha_2(i)^{-1}\mu_0\alpha_1(i)$ centralise $h_i$
% archive sheet number: 7
\scanpage{9}
\diagram{\scriptsize\begin{tikzcd}[ampersand replacement=\&,column sep=.8em]
0 \arrow[r] \& \deleted{\unsure{\mathbb Z(G)}}\,0 \arrow[r] \& \mathbb Z^I \arrow[r,dashed] \&
\underset{\text{Teichmüller spécial}}{TS(C,D,\varphi)} \arrow[r,two heads,"\psi"] \&
\underset{\substack{\text{Teichmüller}\\\text{à trous}}}{T(C,D,\varphi)} \arrow[r] \&
\underset{\substack{\text{groupe}\\\text{symétrique}}}{\unsure{\mathfrak S_I|\mathfrak G_I}} \arrow[r] \& 1
\end{tikzcd}}
% slanted
\margin{nul sauf pour la sphère à 1 ou 2 trous}
Les $(v,\alpha)\in TS$ dans le noyau de $\psi:TS\to T$, i.e. tels que $\widetilde v=\mathrm{id}$, sont représentables par des $(v,\alpha)$ avec $v=\mathrm{id}$, donc
\[
h_i=\mathrm{int}\,(\alpha(i))\,h_i
\]
Sauf dans le cas de la sphère à $1$ trou, \deleted{\unsure{on|on a}} le centralisateur de $h_i$ est le groupe qu'il engendre \uline{et} ce groupe est libre de rang $1$, donc on aura alors $\deleted{\unsure{\in}}$
\[
\alpha(i)=h_i^{\lambda_i}\qquad\lambda_i\in\mathbb Z\ \text{bien déterminé}
\]
On aura
\[
(\mathrm{id}_G,(h_i^{\lambda_i}))\sim(\mathrm{id}_G,(h_i^{\lambda'_i}))
\]
sss $\exists\mu\in Z(G)$ tel que
\[
 h_i^{\lambda'_i}=\mu\,h_i^{\lambda_i}
\]
Dans tous les cas sauf la sphère à $2$ trous (cylindre) $Z(G)=0$, donc la condition précédente est $(\lambda_i)=(\lambda'_i)$, \inserted{i.e. $\mathrm{Ker}\,\psi\simeq\mathbb Z^I$} Dans le cas de la sphère à $2$ trous, $Z(G)=G\underset{?}{\simeq}\mathbb Z$ \deleted{la condition est que} \inserted{$(h_1,h_2)=(+1,-1)$}
\[
\deleted{\mu=\unsure{h_i|\ell_i}\quad h_i^{\lambda'_i}\,/\,h_i^{\lambda_i}\quad\unsure{\text{i.e. de}}}
\qquad\deleted{\unsure{\text{or}}\quad h_1^{\lambda_1},h_2^{\lambda_2}}
\]
les conditions sont alors (notation additive de $\mathbb Z$)
\[
\begin{aligned}
\lambda'_1&=\mu+\lambda_1\\
-\lambda'_2&=\mu-\lambda_2\qquad\text{i.e.}\quad\lambda'_2=-\mu+\lambda_2
\end{aligned}
\]
i.e. $\mathrm{Ker}\,\deleted{\unsure{\psi}}\,(\mathbb Z^I\to\mathrm{Ker}\,\psi)\simeq\mathrm{Ker}\,(\mathbb Z^I\xrightarrow[\text{opération somme}]{}\mathbb Z)$

donc \qquad $\mathrm{Ker}\,\psi\simeq\mathbb Z$ \quad (canoniquement)

Enfin dans le cas de la sphère à $1$ trou, $G=0$ donc $\alpha=0$, donc $TS(C,D,\varphi)=0$

On trouve d'autre part, trivialement,
\[
\pi_1(\underline{\mathrm{End}}_{D\text{ fixé}}(C,D,\varphi))=1
\]
% archive sheet number: 8
\scanpage{11}
% margin note written sideways
\margin{Catégorie ou \unsure{isotopique|multigroupe|topologique} des tubes \unsure{orientés|orientables|orientées}}
\(\left\{\vcenter{\hbox{\begin{minipage}{.89\linewidth}
\textcircled{a} $(\Sigma,\omega)$ $1$-complexe combinatoire \inserted{(fini)}, circularisé ``ordinaire'' (ss pts isolés ni boucles isolés)

\textcircled{c} $\Pi_d$ $\mathbb Z$-groupoïde \deleted{\unsure{(à $\pi_0(\Pi_b)$ fini)|(à $\pi_0(\Pi_d)$ fini)}} $1$-spécial

\textcircled{b} $\Pi_b$ groupoïde $1$-spécial (i.e. $\pi_0(\Pi_b)$ fini, les groupes fond. isom à $\mathbb Z$)


\end{minipage}}}\right.\)

\textcircled{a$'$} \(\left\{\vcenter{\hbox{\begin{minipage}{.89\linewidth}
$\deleted{\unsure{\Pi}}$ $\partial_t(\Sigma,\omega)$ $1$-complexe combinatoire \deleted{\unsure{circul.}} \uline{orienté} dit ``des co-circuits'' de $\Sigma$.
\[
 \Pi(\partial_t(\Sigma,\omega))\qquad \mathbb Z\text{-groupoïde }1\text{-spécial}
 \quad\xrightarrow{\varphi_0}\Pi(\Sigma)
\]


\end{minipage}}}\right.\)

\textcircled{b$'$}
\(\left\{\vcenter{\hbox{\begin{minipage}{.89\linewidth}

\[
 \overset{\deleted{\unsure{\sim}}}{\Pi}_b=\coprod_{i\in B}
 \underset{\text{comp. connexes}}{\underline{\Pi_b(i)}}
 \qquad \Lambda(i)=\pi_1\Pi_b(i)
\]
\hfill (groupe isom. à $\mathbb Z$)
\[
 \widetilde\Pi_b(i)=\deleted{\unsure{\mathbb Z}}
 \coprod_{\alpha\in\mathrm{Isom}_{\rm gr}(\mathbb Z,\Lambda_{\deleted{\unsure{b}}}(i))}
 \Pi_b(i)
\]
$\mathbb Z$-groupoïde \deleted{\unsure{connexe|grossier}} \inserted{à deux} composantes connexes
\diagram{\begin{tikzcd}[ampersand replacement=\&]
\widetilde\Pi_b \arrow[d,"\varphi_b"] \arrow[r,phantom,"="] \&
\displaystyle\coprod_{i\in B}\widetilde\Pi_b(i) \& \mathbb Z\text{-groupoïde }1\text{-spécial}\\
\Pi_b \& \&
\end{tikzcd}}


\end{minipage}}}\right.\)

\textcircled{c$'$}
$\mathrm{Cat}\,\pi_0(\Pi_d))\xleftarrow{\varphi_c}\Pi_d$

\[
 T(\Sigma,\omega,\Pi_d,\Pi_b)=(C,\underset{\nearrow\,\mathbb Z\text{-groupoïde }1\text{-spécial}}{D},\varphi)
\]
\[
\begin{aligned}
D&=\Pi(\partial_t(\Sigma,\omega))\amalg\widetilde\Pi_b\amalg\Pi_d\\
C&=\Pi(\Sigma)\amalg\Pi_b\amalg\mathrm{Cat}\,(\pi_0(\Pi_d))\\
\varphi&=\varphi_0\amalg\varphi_b\amalg\varphi_d
\end{aligned}
\]
\margin{\uline{NB} $\varphi$ est foncteur bordant}

\noindent\rule{\linewidth}{.4pt}

\textcircled{d} $C'\unsure{,}$ groupoïde

\textcircled{e} $\varphi':D\longrightarrow C'$ \qquad $(\varphi'=(\varphi'_0,\varphi'_b,\varphi'_d))$

\hspace{3em}foncteur bordant
% archive sheet number: 10
\scanpage{12}
Catégorie ou \unsure{isotopique|multiplique|multigroupe} des $1$-complexes topologiques compacts \& subdivisions\ldots.

\[
\left\{\begin{aligned}
&K_{\deleted{\unsure{0}}1}\quad\text{$1$-complexe \inserted{combinatoire,} sans boucles isolés (``\unsure{source|sans}'' défini sur \unsure{$H$-orits|$\frac12$-arêtes} \ldots)}\\
&\Pi_b\quad\deleted{\unsure{\mathbb Z}}\ \text{groupoïde $1$-spécial}
\end{aligned}\right.
\]
\[
\underline{\mathrm{Hom}}\big((K_1,\Pi_b),(K'_1,\Pi'_b)\big)
=\mathrm{Cat}\,(\mathrm{Isom}(K_1,K'_1))\times
\underline{\mathrm{Equ}}_{\deleted{\unsure{\mathbb Z}}}(\Pi_b,\Pi'_b)
\]
% archive sheet number: 11
\scanpage{14}
1) Soit une extension de groupes topologiques
\[
1\to N\to G\to H\to1 \tag*{(E)}
\]
et un hom de groupes topologiques
\[
 H'\xrightarrow{u}H
\]
d'où par image inverse une extension \inserted{$G'$} de $H'$ par $N$ et un hom d'extensions
\diagram{\begin{tikzcd}[ampersand replacement=\&,column sep=1.25em,row sep=1.5em]
(E) \& 1 \arrow[r] \& N \arrow[r] \& G \arrow[r] \& H \arrow[r] \& 1\\
(E') \& 1 \arrow[r] \& N^{\deleted{\unsure{\prime}}} \arrow[u,"\mathrm{id}"'] \arrow[r] \& G' \arrow[u] \arrow[r] \& H' \arrow[u,"u"'] \arrow[r] \& 1
\end{tikzcd}}
Passant aux $\pi_0$, on en déduit un hom de suites exactes \deleted{(sans $1$ à gauche}
\diagram{\begin{tikzcd}[ampersand replacement=\&,column sep=1.25em,row sep=1.5em]
(S) \& \pi_1 H \arrow[r] \& \pi_0N \arrow[r] \& \pi_0G \arrow[r] \& \pi_0H \arrow[r] \& 1\\
(S') \& \pi_1 H' \arrow[u] \arrow[r] \& \pi_0N \arrow[u,equal] \arrow[r] \& \pi_0G' \arrow[u] \arrow[r] \& \pi_0H' \arrow[u] \arrow[r] \& 1
\end{tikzcd}}
dont le carré de droite donne lieu à un hom
\[
\underset{\substack{\Vert\\\pi_0(G\times_H H')}}{\pi_0G'}\xrightarrow{\unsure{\varepsilon|c}}
\pi_0G\times_{\pi_0H}\pi_0H'.
\]

Il est facile de voir sur des exemples que cet hom (qui est tjrs surjectif) n'est pas néc. bijectif. Ainsi, la seule connaissance de \inserted{(disons)} la suite $(S)$ et des $\pi_i(H')\to\pi_i(H)$ ne permet pas, en général, de \deleted{\unsure{calculer}} déterminer $\pi_0(G')$. Notre propos est de préciser des données algébriques simples, \unsure{rel$^{\rm t}$es|relatives} à l'extension $(E)$, qui permettent de \deleted{calculer} \inserted{déterminer} \inserted{(mod isom)} $\pi_0(G')$ connaissant $\pi_0\unsure{(H)|H}\to\pi_0H'$ et $\pi_1H\to\pi_1H'$ \unsure{induits|induit} par $u$.
% archive sheet number: 13
\scanpage{15}
$\alpha$) Réduction au cas $N$ discret.

Observons que $N^0$ est invariant dans $G$ (et dans $G'$) et que
\[
\left\{\begin{aligned}
\pi_0 G&=\pi_0 G/N^0\\
\pi_0 G'&=\pi_0 G'/N^0
\end{aligned}\right.
\]
; d'ailleurs $G'/N^0$ comme extension de $H$ par $N/N^0=\pi_0N$ l'image inverse de $G/N^0$ (extension de $H$ par $\pi_0N$) via $H'\to H$. D'ailleurs le diagramme commutatif
\diagram{\begin{tikzcd}[ampersand replacement=\&]
 H \arrow[r] \& H/N^0\\
 H' \arrow[u,"u"] \arrow[r] \& H'/N^0 \arrow[u,"\widehat u"']
\end{tikzcd}}
donne un carré comm.
\diagram{\begin{tikzcd}[ampersand replacement=\&]
 \pi_i H \arrow[r] \& \pi_i(H/N^0)\\
 \pi_i H' \arrow[u,"\pi_i u"] \arrow[r] \& \pi_i(H'/N^0) \arrow[u,"\pi_i\widehat u"']
\end{tikzcd}}
donc la donnée de $\pi_0u$ équivaut à celle de $\pi_0\widehat u$ (car les $\to$ horiz. sont des isom) et celle de $\pi_1u$ \uline{implique} celle de $\pi_1\widehat u$ (car les $\to$ hor. sont des épi) \deleted{Remplaç} Divisant tout par $N^0$, on est donc ramené au cas où $N$ est \uline{discret}. Dans ce cas la suite $(S)$ revient à \deleted{\unsure{des}} isom
\[
\left\{\begin{aligned}
&\pi_iG\simeq\pi_iH\qquad i\geq2\\
&\pi_1G=\mathrm{Ker}\,(\pi_1 H\to\underset{\substack{\Vert\\N}}{\pi_0N})\\
&\pi_0G)\ \text{est ext. de }\pi_0H\ \text{par }\mathrm{Coker}\,(\pi_1 H\to N)
\end{aligned}\right.
\]
et de $=$ pour $\deleted{\unsure{G}}$ les $\pi_iG'$. \deleted{Notre propos} Supposons connus \deleted{\unsure{les}} $\pi_i H'$, on trouve les $\pi_i G$ si $i\geq2$, et $\pi_1G'$ comme $\mathrm{Ker}\,(\pi_1H'\to N)$, en notant que $\deleted{\pi_1H'\to}$ l'on a comm. de
\diagram{\begin{tikzcd}[ampersand replacement=\&]
\pi_1 H \arrow[r] \& N\\
\pi_1 H' \arrow[u,"\pi_1(u)"] \arrow[r] \& N \arrow[u,equal]
\end{tikzcd}}
donc $\pi_1H'\to N$ est connu (pour tt $H'\xrightarrow{u}H$) si $\pi_1H\to N$ est connu, en termes du seul $\pi_1u$. Notre propos est \deleted{donc de} \inserted{en} \deleted{\unsure{fait}} de déterminer (pour tt $H'$ et $u:H'\to H$) l'extension
% archive sheet number: 14
\scanpage{16}
$\pi_0G'$ de $\pi_0H'$ par $R'=\mathrm{Coker}\,(\pi_1 H'\to N)$. On fera \unsure{$\cdot$} attention que $\pi_0G$ est lui-$=$ une extension de \deleted{Coker} $\pi_0H$ par $\mathrm{Coker}\,(\pi_1H\to N)$, son image inverse \inserted{$\pi_0 G\times_{\pi_0 H}\pi_0 H'$} par $\pi_0 H'\to\pi_0H$ est une extension de $\pi_0H'$ par $R=\mathrm{Coker}\,(\pi_1 H\to N)$ \deleted{\unsure{$\simeq$ qui}}, non par $R'$. $R$ est un quotient de $R'$ par
\[
\frac{\mathrm{Im}\,(\pi_1 H\xrightarrow{\varphi}N)}
     {\mathrm{Im}\,(\pi_1 H'\xrightarrow{\varphi'}N)}
\simeq \pi_1H/(\mathrm{Ker}\,\varphi+\mathrm{Im}\,\pi_1(u))
\]
donc on a
\[
1\to\pi_1H/(\mathrm{Ker}\,\varphi+\mathrm{Im}\,\pi_1(u))\to\pi_0G'
\xrightarrow{c}\pi_0G\times_{\pi_0H}\pi_0H'\to1
\]
Cel \deleted{\unsure{montre}} \inserted{met} en évidence de façon précise comment le noyau de \unsure{$a$|$c$} est lié au $\pi_1H$ et $\pi_1u$, c'est un quotient de $\mathrm{Coker}\,\pi_1u$.

$\beta$) \uline{Cas $H$ connexe}. Alors l'ext. $G$ est caractérisée, à isom unique près, par un hom
\[
\deleted{\pi\longrightarrow}\qquad
\pi\overset{\mathrm{def}}{=}\pi_1(H)\xrightarrow{\unsure{\rho|\varphi}}\mathfrak z(N)
\]
et peut se décrire à l'aide du groupe \deleted{\unsure{rev}} universel $\widetilde H$ de $H$ (ext. de $H$ par $\pi$) en prenant son ``image directe'' par $\pi\xrightarrow{\varphi}N$

Alors l'extension \deleted{\unsure{de}} $G'$ de $H'$ par $N$ peut s'obtenir en prenant l'image inverse $H'\times_H\widetilde H$ de \deleted{\unsure{l'}}l'ext. $\widetilde H$ de $H$ par $\pi$ via $H'\to H$, et prenant l'image par $\pi\xrightarrow{\varphi}\mathfrak z(N)\hookrightarrow N$. Il n'y a de pb que si $H'$ \deleted{\unsure{est}} n'est pas connexe \inserted{(\unsure{disons|sinon|décrivons} les ext. de $\pi_0H'=1$ sont évidentes !)} -- p. ex $N$ discret. \deleted{discret (plus gén \unsure{si $G$},} \deleted{Il faut que \unsure{de l'on|prendre l'}}
% archive sheet number: 15
\scanpage{17}
\mspage{A}
\[
\begin{aligned}
(a)&\quad -\ 1\to N\to G\to H\to1
 &&\text{ext. de groupes top.}\quad [\,N\text{ discret}\,]\\
(b)&\quad \deleted{\unsure{\varphi}}\quad 1\to\pi\to\widetilde H\to H\to1
 &&\text{ext. de groupes top.}\quad [\,\pi\text{ \rlap{\(\big/\)}discret}\,]\\
(c)&\quad 1\to\widetilde H_0\to\widetilde H\to\unsure{\mathfrak S|\mathfrak G}\to1
 &&\text{ext. de gr. top.}\quad [\,\mathfrak S\text{ discret}\,]
\end{aligned}
\]

On suppose donné (b) et (c) \inserted{et $N$}, on s'intéresse aux ext. $G$ de $H$ par $N$, \uline{munies} d'une trivialisation \inserted{\uline{stricte} (\unsure{isom.|triv.|i.e. comm.} avec pdt. pr. direct)} de l'image inverse de cette extension par $\widetilde H_0\to H$ (composé $\widetilde H_0\hookrightarrow\widetilde H\to H$) i.e. d'un hom $\widetilde H_0\to G$ qui remonte $\widetilde H_0\to H$ et centralise $N$

(\uline{NB} $N$ étant discret, si $\widetilde H_0$ est connexe et simplement connexe, il existe pour $G$ donné une unique telle trivialisation\ldots)

\inserted{On suppose de plus que ce splitting est invariant par les opérations de $\widetilde G$ sur $\widetilde H_0$ (via $\widetilde H$ \unsure{qui} $\deleted{\unsure{\to H}}$) et sur $G$ (via $\widetilde G\to G$) (cf ci-dessous pour déf. de $\widetilde G$)}
Soit $\widetilde G$ l'image inverse de $G$ par $\widetilde H\to H$
\[
1\to N\to\widetilde G\to\widetilde H\to1
\]
La donnée d'une ext $G$ est équivalente à celle de $\widetilde G$, \uline{munie} d'une trivialisation \inserted{stricte} sur $\pi$. La \deleted{trivialisation de} structure supplémentaire sur l'ext. $G$ (trivialisation de l'image inverse par $\widetilde H_0\to H$) équivaut alors à une trivialisation \inserted{stricte} de $\widetilde G$ sur le sous-groupe $\widetilde H_0$. Donc au résumé on doit étudier les ext. $\widetilde G$ de $\widetilde H$ par $N$, trivialisées \inserted{strictemt} simultanément sur $\pi$ \uline{et} sur $\widetilde H_0$ [rappelons nous que cette dernière structure, \deleted{\unsure{triv.}} \unsure{fixant|spécifique|sur} $\widetilde H_0$, est automatique dans les conditions dites]. Or une ext. \deleted{\unsure{trivial.}} \inserted{$\widetilde G$ de} $\widetilde H$ par $N$ trivialisée \inserted{\unsure{strit|strict|strictement}} sur $\widetilde H_0$ équivaut à une ext. de $\mathfrak S$ par $N$, soit $\widetilde G_0$ :
\[
1\to N\to\overset{\deleted{\unsure{\sim}}}{G}_0\to\mathfrak S\to1
\]
($\widetilde G$ est l'image inverse par $\widetilde H\to\mathfrak S$). De plus, il faut donner un splittage \inserted{strict} du diagramme
% archive sheet number: 16
\scanpage{18}
\diagram{\begin{tikzcd}[ampersand replacement=\&,column sep=small,row sep=small]
1 \arrow[r] \& N \arrow[r] \& \overset{\deleted{\unsure{\sim}}}{G}_0 \arrow[r] \& \unsure{\mathfrak S|\mathfrak G} \arrow[r] \& 1\\
\& \& \& \pi \arrow[u] \arrow[ul,dashed,"\varphi"'] \&
\end{tikzcd}}
invariant par $\widetilde G$ [opérant sur $\pi$ (via $\widetilde H$ et $\widetilde G\to\widetilde H$) et \deleted{\unsure{$N$}} \inserted{$G_0$} (via $G$ \inserted{et $\widetilde G\to G_0$})]
% slanted
\margin{si $\pi$ \unsure{comm.|connexe} et discret, l'opération \inserted{de $\unsure{G|\widetilde G}$} sur $\pi$ se fait via $\unsure{H|\widetilde H}$, donc via $\pi_0(H)$ (ou $\pi_0(\widetilde H)$) et celle de $\unsure{G|\widetilde G|\mathfrak G}$ sur $\mathfrak z(N)$ se fait via $\mathfrak S$ (\unsure{car} $N$ \deleted{\unsure{agit}} \unsure{opère} triv$^{\rm t}$) et si $\mathfrak S$ discret et $\widetilde H_0$ connexe, on a $\pi_0(H)=\pi_0(\widetilde H)=\mathfrak S$) donc la condition est que $\varphi:\pi\to\mathfrak z(N)$ soit $\mathfrak S$-hom.}

Cas particulier : $\pi\subset\widetilde H_0$, i.e. $\pi\to\mathfrak S$ est trivial, alors la donnée de $\varphi$ équivaut à la donnée d'un hom
\[
\varphi:\pi\to\mathfrak z(N),
\]
\noindent\rule{\linewidth}{.4pt}

\diagram{\scriptsize\begin{tikzcd}[ampersand replacement=\&,column sep=small,row sep=small,/tikz/execute at end picture={\draw[thin] (current bounding box.north west) -- (current bounding box.south east);}]
\& \& \& \deleted{\widetilde H_0} \arrow[dl,hook] \arrow[dr] \& \&\\
1 \arrow[r] \& \pi \arrow[r] \& \widetilde H \arrow[rr] \& \& H \arrow[r] \& 1\\
1 \arrow[r] \& \pi \arrow[u,equal] \arrow[r] \& \widetilde H' \arrow[u,"\widetilde f"] \arrow[rr] \& \& H' \arrow[u,"\deleted{\unsure{f}}", "f"'] \arrow[r] \& 1\\
\& \& \& \widetilde H'_0=f^{-1}(\widetilde H_0) \arrow[uuu] \arrow[ul,hook] \arrow[ur] \& \&
\end{tikzcd}}
\[
\deleted{\unsure{\mathrm{int}(\widetilde g)}\,\varphi(\pi)=\unsure{-}}
\]
\diagram{\scriptsize\begin{tikzcd}[ampersand replacement=\&,row sep=small,column sep=small,/tikz/execute at end picture={\draw[thin] (current bounding box.north west) -- (current bounding box.south east);}]
\& \& 1 \& \&\\
1 \arrow[r] \& \pi \arrow[r] \arrow[dr,"\varphi"'] \& \deleted{\widetilde H} \arrow[u] \arrow[r] \& H \arrow[r] \& 1\\
1 \arrow[r] \& \deleted{\unsure{\pi}} \arrow[u] \arrow[r] \arrow[d] \& \widetilde G \arrow[u] \arrow[r] \& \unsure{G?|G} \arrow[u] \arrow[r] \& 1\\
\& \deleted{\unsure{N=}} \& N \arrow[u] \arrow[r] \& N \arrow[u] \&\\
\& \& 1 \arrow[u] \& \&
\end{tikzcd}}

\diagram{\scriptsize\begin{tikzcd}[ampersand replacement=\&,row sep=small,column sep=small,/tikz/execute at end picture={\draw[thin] (current bounding box.north west) -- (current bounding box.south east);}]
1 \arrow[r] \& \mathscr H_1/\mathscr H_1\cap\mathscr H_2 \arrow[r] \& \overset{\unsure{p_2|j_2}}{\mathscr H/\mathscr H_2} \arrow[r] \& \overset{\unsure{p_1|j_1}}{\mathscr H/\mathscr H_1\mathscr H_2} \arrow[r] \& 1\\
1 \arrow[r] \& \mathscr H_1 \arrow[r,hook] \& \deleted{\mathscr H} \arrow[u] \arrow[r] \& \mathscr H/\mathscr H_1 \arrow[u] \arrow[r] \& 1\\
1 \arrow[r] \& \mathscr H_1\cap\mathscr H_2 \arrow[u] \arrow[r] \& \mathscr H_2 \arrow[u,hook] \arrow[r] \& \mathscr H_2/\mathscr H_1\cap\mathscr H_2 \arrow[u] \arrow[r] \& 1\\
\& 1 \arrow[u] \& 1 \arrow[u] \& 1 \arrow[u] \&
\end{tikzcd}}

\diagram{\scriptsize\begin{tikzcd}[ampersand replacement=\&,row sep=tiny,column sep=small,/tikz/execute at end picture={\draw[thin] (current bounding box.north west) -- (current bounding box.south east);}]
\& \& 1 \& 1 \&\\
\& \& \mathfrak S \arrow[u] \arrow[r,equal] \& \mathfrak S \arrow[u] \&\\
1 \arrow[r] \& \pi \arrow[r] \& \deleted{\widetilde H} \arrow[u] \arrow[r] \& H \arrow[u] \arrow[r] \& 1\\
1 \arrow[r] \& \pi \arrow[u,equal] \arrow[r] \& \widetilde H_0 \arrow[u] \arrow[r] \& H_0 \arrow[u] \arrow[r] \& 1\\
\& \& 1 \arrow[u] \& 1 \arrow[u] \&
\end{tikzcd}}
\diagram{\scriptsize\begin{tikzcd}[ampersand replacement=\&,column sep=1em]
\mathrm{Ext}(\mathfrak S,N) \arrow[d] \&[-2em] {}\times\mathrm{Hom}_{\mathfrak S}(\pi,\mathfrak z(N)) \arrow[r,"\sim"] \& \deleted{\mathrm{Ext}(H,\mathrm{mod}\,\widetilde H_0,N)} \arrow[d]\\
\substack{\text{Esp. homogène}\\\text{sur }H^2(\mathfrak S,N)} \& \& \mathrm{Ext}(H,N)\\
\& \& \substack{\text{espace homogène}\\\text{sur }H^2(H,\mathfrak z(N))}
\end{tikzcd}}
\diagram{\scriptsize\begin{tikzcd}[ampersand replacement=\&,row sep=tiny,column sep=tiny]
\& \& \deleted{\unsure{\mathscr H}} \& \cdot \&\\
1 \arrow[r] \& N \arrow[r] \& G \arrow[r] \arrow[d] \& \widetilde H \arrow[u] \arrow[r] \& 1\\
\& \& H \& \overset{\deleted{\unsure{\sim}}}{H} \&
\end{tikzcd}}
\[
\deleted{H^2(\mathfrak S,\mathfrak z N)\to H^{2}(H,\mathfrak z N)}
\]
car $H\to\mathfrak S$
\[
\deleted{\mathrm{Ext}(\mathfrak S,N)\times\mathrm{Hom}(\pi,N)\to H^2(H,\mathfrak z(N))}.
\]
% archive sheet number: 17
